\section{二つの寮生活を経て――熊野寮一年目の雑感}
\rightline{（文責:イカリング）}

\subsection{通算4年目の寮生活}
　早いものでもう熊野寮に入ってから1年近く経とうとしている。今年のはじめは高校生だったことが驚きだ。実は私は高校生の間も寮生活をしていたので通算4年目の寮生活になる。今のところ人生の四分の一ほどを寮生活しているのだ。今後もあまり寮を出るつもりはないので修士号を最短で取ったとしたら寮生活が通算9年になるそうだ(わおびっくり)。そんな熊野寮について、前の寮とはいくつか違う点があって面白いのでここに書き綴ろうと思った。

\subsection{男女混合寮という驚き}
　はじめに、熊野寮は男女混合寮であることだ。え？そんなこと？と思った方もいるかもしれない。確かに熊野寮は、ぼろいし、自治寮だし、汚いし、ぼろくて汚いなどの特徴があるが、男女混合というのは私にとってかなり大きな要素であった。私の高校は男子校で、女子部屋や女子寮はおろかそもそも女子が存在しなかった。そんな中で暮らしていたので世の中には男子寮と女子寮しか存在しないと思っていたのだ。そんなものだったので熊野寮の入寮初日に女子部屋の存在を知ってびっくりした。こういうところにいると当たり前のように感じてしまうのだが、京大生が集まっているからこそできることなのだと思う。こんなことを言うと自称リベラリスト(笑)の方々に怒られるかもしれないが、世の中には共学の「学校」でさえ男女仲が険悪になるところも多くあるのを忘れてはならない。ましては完全な生活空間である「寮」で男女混合を実現しようというのは決して簡単なことではなかったのだと思う。その点ではリベラリズム的規範意識が比較的高い京大生が集まるという土壌がよかったのだろう。男女混合寮特有の空気感の中でのびのびと生活をするのも男子寮でのびのびと生活していた高校寮のときとはまた違った楽しさがある。男女混合寮というのは今の熊野寮特有の空気感が得られるいい要素なのだと思う。

\subsection{自治寮という仕組み}
　次に、自治寮であることだ。高校寮で自治なんていうのはありえなかった。高校寮では寮教諭の方が10人ほどいて、夜廊下を走り回ったりしたら怒られたりしていた。教育なのか虐待なのかあいまいなルールもあって、先生方はそのルールに則ったり則らなかったりして寮内の秩序を守ってくださった。実際、これによって大きなトラブルは速やかに解決したし、大事にならずに済んだのでいい面もあるのだと思う。自治ができるというのは素晴らしいことなのだが、内側の問題を解決するのに時間がかかるというのは致命的な欠点であるように思う。そうしたトラブルもまた生活の一部であるというようにして受け入れることも大事であろう。負の側面を書いてしまったが、自治寮であることならではの良さはある。なんといってもルールを自分たちで作ることができることであろう。現場を知らない人がルールを決めるというのは往々にしてあることなのだが、それでは当然現場に合わないルールが出来上がってしまう。こういうのはいくつでも心当たりがあると思われる。文科省の考えることとか全部そうですよね(共通テストをやって思考力が養われるわけないし...)。自分たちでルールを決められるということはそのような不思議な縛りがないということである。今後熊野寮は変わっていくのだろうがこの良さは今後にも残してもらいたいと思う。

\subsection{学年混合のつながり}
　また、学年混合であることも目新しかった。高校寮では入学と同時に同じ階(今は大部屋)に押し込んで学校で顔を見る同級生と生活をしていた。学年が上がるタイミングで部屋移動が行われ、4人でグループを作って一年を過ごすというサイクルで生活をしていた。上級生、下級生と顔を合わせる機会が熊野寮でのそれよりもはるかに少なかったのだ。熊野寮に入ってからはむしろ学年の違う方と関わることが非常に増えた。今の私は同学年の友達と顔を合わせる機会はたまに行くサークルと授業の教室、寮のコンパくらいしかない。横のつながりよりも縦のつながりの方が増えたのは高校寮とは全く異なる点である。

\subsection{入寮を考えている君へ}
　ここまで書いたが今回はこのくらいにしようと思う。最後に、これを読んでいる君はきっと熊野寮への入寮を考えているのだと思うので君に伝えたいことを伝えて終ろうと思う。君が合格して熊野寮に入ることになったとして、君は熊野寮についておそらくよく知らないまま入ることになると思う。でもそれでいいんだ。もし興味があるならぜひ入寮して私たちのことを知ってほしい。君と会える日を楽しみにしている。A2談話室においで。歓迎するよ！